
\lesson{16/11}
So the general strategy is to write down a Lagrangian, define its Feynman rules
and use them to write down transition amplitudes.
\chapter{Path Integrals for fermionic fields.}
The fields for fermions are \textit{Dirac Spinors} $\psi$ and for the canonical
quantization the commutation relations were changed to anticommutators: how can
this be done with path integrals?

\section{Grassman Variables}
Grassman numbers are \textit{anticommuting numbers},
\[
    \left\{\vartheta,\vartheta \right\}=0\iff \vartheta^2 = 0;
\]
They are also called G-numbers, opposed to C-numbers which are complex.

Functions of Grassman variables are thus simple to expand in series: since
$\vartheta^2=0$, $\vartheta^{n>2}=0$ and the series expansion is
\[
    f(\vartheta ) = c +\alpha \vartheta
\]
\textbf{Note} If $f$ is a G-number (C-number), $c$ is a G-number (C-number) and
$\alpha $ is a C-number (G-number).

Since $\dv[]{\vartheta }{\vartheta }=1$, the derivative operators anticommute
\[
    \left\{\dv[]{}{\vartheta },\dv[]{}{\vartheta }\right\}=0.
\]
A desirable property is that the integral over a G-variable is
\textit{invariant under translation}
\[
    \int_{}^{} \dd \vartheta f(\vartheta +\eta ) = \int_{}^{} \dd \vartheta
    f(\vartheta )
\]
This implies that
\[
    \int_{}^{} \dd \vartheta (c+\alpha \vartheta ) = \int_{}^{} \dd \vartheta 
    (c+\alpha \vartheta )+\int_{}^{} \dd \vartheta \alpha \eta \iff
    \int_{}^{} \dd \vartheta =0.
\]
The integration measure can be defined so that
\[
    \int_{}^{} \dd \vartheta \vartheta  = 1.
\]
Moreover, different G-variables anticommute
\begin{gather*}
    \left\{\vartheta_i,\vartheta_j\right\}=0\\
    \left\{\dv[]{}{\vartheta _i},\dv[]{}{\vartheta _j}\right\}=0
\end{gather*}
so the Taylor expansion of the function of $n$ G-variables is
\[
    f(\vartheta_i) = c+ \sum_{i}^{}\alpha_i\vartheta _i + \sum_{ij}^{} \beta _{ij}
    \vartheta _i\vartheta _j + \sum_{ijk}^{} ...+... +
    \gamma \vartheta_1\vartheta_2...\vartheta_n
\]
since higher order terms are all null.

\textbf{Note} The derivative and integral operators \textit{must act directly
on the variable} and commute with G-variables, i.e.
\[
    \dv[]{}{\vartheta_1}\vartheta_2\vartheta_1 = -\vartheta_2\dv[]
    {\vartheta _1}{\vartheta _1} = -\vartheta_2
\]
\[
    \int_{}^{} \dd \vartheta_1\dd\vartheta_2 \vartheta_1\vartheta_2 = 
    -\int_{}^{} \dd\vartheta_1\vartheta_1\int_{}^{} \dd \vartheta_2\vartheta_2 =-1
\]
\subsection{Gaussian integrals}
Since the function $e^{-\vartheta_2}$ would be the constant $1$, consider
\[
    I=\int_{}^{} \dd \vartheta_1\vartheta_2 e^{-\vartheta^TM\vartheta}
\]
where
\[
    \vartheta = \begin{pmatrix*}
    \vartheta_1\\\vartheta_2
    \end{pmatrix*}
\]
For regular numbers, it can be shown that\footnote{See the section "Determinanti
    funzionali" in the QM2 course in this repository}:
\[
    I_C = \int_{}^{} \dd^2x e^{-x^TMx} = \frac{\pi}{\sqrt{\det(M)}}
\]
Explicitely,
\[
    \vartheta^TM\vartheta = \vartheta_1(m_{11}\vartheta _1+m_{12}\vartheta_2)+
    \vartheta_2(m_{21}\vartheta_1+m_{22}\vartheta_2)=(m_{12}-m_{21})\vartheta_1
    \vartheta_2
\]
so the only part of $M$ that is important is the \textit{antisymmetric part}.
If $M$ is antisymmetric then
\[
    \vartheta^TM\vartheta = 2m_{12}\vartheta_1\vartheta_2.
\]
Expanding the integral yields
\[
    I=\int_{}^{} \dd\vartheta_1\dd\vartheta_2 (1-\vartheta^TM\vartheta)=
    -2m_{12}\int_{}^{} \dd\vartheta_1\dd\vartheta_2\vartheta_1\vartheta_2=2m_{12}=
    2\sqrt{\det(M)}
\]
If the G-variables are odd, then the integral is null; for all even numbers
$2n$ of G-variables, considering the \textit{antisymmetric part of} $M$ the
general result is
\[
    I_{2n}=\int_{}^{} \dd\vartheta_1...\dd\vartheta_{2n} e^{-\vartheta^TM\vartheta} = 
    \int_{}^{} \dd \vartheta_1...\dd\vartheta_{2n} \exp\left[ \sum_{\alpha=1}^{n}
    \vartheta_\alpha^TM_\alpha \vartheta_\alpha\right]
\]
where\[
    M_\alpha  = \begin{pmatrix}
        0&\lambda_\alpha\\
        -\lambda_\alpha&0
    \end{pmatrix}
\] 
are the blocks of the Jordan form of $M$ and $\vartheta_\alpha =
(\vartheta_\alpha\quad\vartheta_{\alpha+1})$ (credo).
The integral can be factorized and expanded in
\[
    \text{[buchino]},
\] yielding
\[
    I_{2n} =2^{2n}\sqrt{\det(M)}
\]
Consider now
\[
    I = \int_{}^{} \dd \vartheta_1\dd\vartheta_{2n}
    \exp\left[ -\vartheta^TM\vartheta+\eta^T\vartheta+\vartheta^T\eta\right]
\]
the exponent can be rewritten as 
\[
    \vartheta'^TM\vartheta' - \eta^TM^{-1}\eta,\;\text{with }
    \vartheta' = \vartheta-M^{-1}\eta
\]
using the translational invariance yields
\[
    I=2^n\sqrt{\det(M)}e^{\eta^TM^{-1}\eta }
\]
\section{Dirac Fields}
The Dirac action is
\[
    S_0 = \int_{}^{} \dd^4x\bar\psi(i\slashed{\pd}-m)\psi 
\]
and the equations of motion are simply
\begin{gather*}
    (i\slashed{\pd}-m)\psi =0\\
    i\slashed{\pd}\bar\psi +m\bar\psi =0
\end{gather*}

The partition function for this free theory is
\[
    Z_0=\int \mathcal{D}\psi\mathcal{D}\bar\psi  e^{iS_0}
\]
The goal is again to find Green's functions of the form
\[
    \int \mathcal{D}\psi\mathcal{D}\bar\psi e^{iS_0}\psi_1\psi_2...\bar\psi_1...
    \bar\psi_n
\]
So again a source term, the spinor $\eta$, is used:
\[
    Z_0[\eta,\bar\eta]=\frac{1}{\mathcal{N}}\int\mathcal{D}\psi\mathcal{D}\bar\psi
    \exp\left[ iS_0 + i \left\langle \bar\eta \psi\right\rangle +i\left\langle 
    \bar\psi \eta\right\rangle  \right]
\]
where $\bar\eta =\eta^\dagger\gamma^0$. As before, this is useful since
\[
    -i\fdv[]{}{\bar\eta_\alpha(x)}Z_0[\eta,\bar\eta] =
    \int \mathcal{D}\psi \mathcal{D}\bar\psi \psi_\alpha(x)\exp\left[ iS_0+i
    \left\langle \bar\eta \psi \right\rangle +i\left\langle\bar\psi\eta
\right\rangle \right]
\]
And similarly applying $i\fdv[]{}{\eta_\alpha(x) }$ brings down a
$\bar\psi_\alpha(x)$, where $\alpha=0,1,2,3$ is the component of the spinor (which
will be omitted from now on).

So to compute
\[
    G=\int \mathcal{D}\psi\mathcal{D}\bar\psi \psi(x_1)\psi(x_2)\psi(x_3)e^{iS_0}
\]
one can use
\[
    (-i)^3\fdv[]{}{\bar\eta_3}\fdv[]{}{\bar\eta_2}\fdv[]{}{\bar\eta_1}
    Z_0[\eta,\bar\eta]
\]
with the \textit{caveat} that
$\pdv[]{}{\bar\eta_i}\psi_j = -\psi_j\pdv[]{}{\bar\eta_i}$, which yields a total
of $3$ negative signs, so
\[
    (-i)^3\fdv[]{}{\bar\eta_3}\fdv[]{}{\bar\eta_2}\fdv[]{}{\bar\eta_1}
    Z_0[\eta,\bar\eta]=-G
\]
Since the functional derivatives anticommute, to get $G$ one can also reverse the
order of derivation. For $n$ $\psi$ fields the general formula is
\[
    \eval{G(1,2...n) = (-i)^n \frac{\delta^n}{\delta\bar\eta_1...
    \delta\bar\eta_n}Z_0[\eta,\bar\eta]}_{\eta_i=0,\bar\eta_j=0}
\]

Now we use F.Ts on the exponent of $Z_0$:
\begin{gather*}
    \int_{}^{} \dd^4x \left[\bar\psi (\slashed{\pd}-m)\psi +\bar\eta \psi+\bar
    \psi \eta \right]=\\
    \int_{}^{} \dd^4x \frac{\dd^4k\dd^4k'}{(2\pi)^8}\left[
    \bar{\tilde\psi}(k')(\slashed{k}-m)\tilde{\psi}(k)+\bar{\tilde\eta}(k')
\tilde\psi(k)+\bar{\tilde\psi}(k')\tilde\eta(k)\right]e^{-i(k+k')x}
\end{gather*}
recognizing a $\delta^4(k'+k)$ yields
\[
    \int_{}^{} \frac{\dd^4k}{(2\pi)^4}\left[(\bar{\tilde\psi }(-k)+\bar{
    \tilde\eta}(-k)(\slashed{k}-m)(-1))(\slashed{k}-m)(\tilde\psi(k)+(\slashed{k}-m)
^{-1}\tilde\eta(k))-\bar{\tilde\eta }(-k)(\slashed{k}-m)\tilde\eta(k)\right]
\]      
This is the same form of the exponent of Gaussian integral of the G-numbers, with
\[
    M=\slashed{k}-m
\]
so one can change variable of integration (with Jacobian $\mathbb{I}$) and write
\[
    Z_0[\eta,\bar\eta] =\frac{1}{\mathcal{N}}\int
    \mathcal{D}\psi\mathcal{D}\bar\psi \exp\left[ iS_0 -i\int_{}^{} \dd^4k
    \bar{\tilde\eta }(-k)(\slashed{k}-m)\tilde\eta(k)
    \right] 
\]
Since $\mathcal{N}=\int \mathcal{D}\psi\mathcal{D}\bar\psi e^{iS_0}$ they
cancel out; antitransforming yields
\[
    Z_0[\eta,\bar\eta] = \exp\left[ -\int_{}^{} \dd^4x\dd^4y
    \bar\eta_\alpha(x) S_{\alpha,\beta}(x-y)\eta_\beta(y)\right]
\]
where
\[
    S(x-y)=i\int_{}^{} \frac{\dd^4k}{(2\pi)^4}e^{-ik(x-y)}\frac{\slashed{k}-m}{
    k^2-m^2+i\varepsilon }.
\]
So finally
\[
    Z_0[\eta,\bar\eta]=\exp\left[ -\left\langle \bar\eta(x)S_{xy}\eta(y)
    \right\rangle_{xy}\right]
\]
And expectation values can be found like
\[
    \eval{\bra{0}T(\psi(x_2)\psi(x_2))\ket{0} = (-i)^2\fdv[]{}{\bar\eta_1}\fdv[]{}
    {\bar\eta_2}Z_0[\eta,\bar\eta]}_{\bar\eta_i=0}=0\footnote{mostralo}
\]
This can also be done with different components of different spinors, like
\[
    \bra{0}T(\psi_\alpha(x)\bar\psi_\beta(y))\ket{0} = (-i)\fdv[]{}
    {\bar\eta_\alpha(x)}\eval{i\fdv[]{}{\eta_\beta(y)} Z_0[\eta,\bar\eta]}_{
    \eta,\bar\eta=0}=[\text{simple}]=S_{\alpha\beta}(x-y)
\]
